%% ─────────────────────────────────────────────────────────────────────────────
%% Lettre de motivation — Talvin Ackbaraly × Valtech
%% Stage : Intelligence Artificielle (Talent Program)
%% Même charte que le CV (XCharter, un seul accent, filets fins)
%% Compile: latexmk -pdf lettre-motivation-valtech.tex
%% ─────────────────────────────────────────────────────────────────────────────
\documentclass[11pt, a4paper]{extarticle}

\usepackage[T1]{fontenc}
\usepackage[utf8]{inputenc}
\usepackage[french]{babel}
\usepackage[top=1.6cm, bottom=1.6cm, left=2.2cm, right=2.2cm]{geometry}
\usepackage{XCharter}
\usepackage[letterspace=80]{microtype}
\usepackage{ragged2e}
\usepackage{xcolor}
\usepackage[hidelinks]{hyperref}

\hypersetup{pdftitle={Lettre de motivation — Talvin Ackbaraly}, pdfauthor={Talvin Ackbaraly}}

%% ── Colour system: identique au CV ──────────────────────────────────────────
\definecolor{accent}{HTML}{2563A0}
\definecolor{ink}{HTML}{1A1A1A}
\definecolor{soft}{HTML}{6E6E6E}
\definecolor{hair}{HTML}{C9C9C9}

\hypersetup{colorlinks=true, urlcolor=accent, linkcolor=accent}
\AtBeginDocument{\color{ink}}

\pagestyle{empty}
\setlength{\parindent}{0pt}
\setlength{\JustifyingParindent}{0pt}
\setlength{\parskip}{8pt}
\linespread{1.08}
\microtypesetup{expansion=false}

\begin{document}

%% ── En-tête : expéditeur ─────────────────────────────────────────────────────
%% Coordonnées seulement : l'école et les liens sont déjà dans le CV
{\color{accent}\fontsize{22}{26}\selectfont\scshape\textls[40]{Talvin Ackbaraly}}\par\vspace{3pt}
{\setlength{\parskip}{0pt}
94270 Le Kremlin-Bicêtre\\
\href{tel:+33769389008}{\textcolor{ink}{07 69 38 90 08}}\\
\href{mailto:talvin.ackbaraly@icloud.com}{\textcolor{ink}{talvin.ackbaraly@icloud.com}}\par}
\vspace{-4pt}{\color{hair}\rule{\textwidth}{0.4pt}}\par

%% ── Destinataire et date ─────────────────────────────────────────────────────
\vspace{2pt}
%% Bloc calé contre la marge droite, texte aligné à gauche dans le bloc
\hfill\begin{tabular}[t]{@{}l@{}}
  \textbf{Valtech}\\
  Équipe Talent Acquisition\\
  Paris
\end{tabular}

Le \today

\vspace{2pt}
\textbf{Objet : candidature au stage Intelligence Artificielle (Talent Program) — 6 mois à partir
de février 2027}

\vspace{2pt}
Madame, Monsieur,

\justifying
Étudiant en dernière année à EPITA, majeure Image, je vous adresse ma candidature pour le stage
Intelligence Artificielle de votre Talent Program. Ce qui m'attire chez Valtech, c'est le chemin
que décrit votre offre : partir d'un cas d'usage métier réel, prototyper une solution d'IA
générative, la mesurer, puis la présenter à un client. C'est un chemin que j'ai déjà parcouru
par morceaux, et que je veux maintenant suivre de bout en bout sur des projets LLM.

Côté IA, mon projet de fin d'études avec la Bibliothèque nationale de France détecte et classe les
illustrations des documents de Gallica, récupérés via l'API IIIF. J'y suis responsable de la
détection : entraînement sur 10\,000 images annotées, comparaison de plusieurs modèles, réglage des
hyperparamètres, jusqu'à 94,7\,\% de mAP50. J'ai aussi implémenté un U-Net from scratch en
PyTorch, sans poids pré-entraînés. Ces projets m'ont appris qu'un modèle ne vaut que par la façon
dont on l'évalue, et c'est ce réflexe que je veux appliquer aux protocoles d'évaluation de
solutions RAG ou agentiques.

Côté développement, je sais transformer un modèle en quelque chose que l'on peut montrer et
utiliser. En freelance, je conçois et j'exploite seul une application de gestion de stock pour un
client professionnel, du recueil des besoins à la mise en production (Next.js, NestJS,
PostgreSQL). Chez Studiocanal, j'ai développé des fonctionnalités full-stack pour une plateforme
utilisée par plus de 100 collaborateurs, et chez Hitachi Solutions, j'ai réalisé un prototype Django
pour Sonepar, premier client du cabinet, de la modélisation UML à la démonstration.

Enfin, je travaille au quotidien avec des agents de code comme Claude Code, Codex ou Antigravity.
Leur usage m'a montré combien la qualité d'un agent dépend de son contexte, de ses outils et des
vérifications qui l'entourent. Je souhaite désormais passer de leur utilisation à leur conception,
et les projets d'agents, de RAG et de workflows orchestrés de votre équipe IA sont exactement le
cadre dans lequel je veux apprendre à le faire, avec la perspective de poursuivre en CDI.

Je suis disponible pour un stage de six mois à Paris à partir de février 2027, et je serais
heureux d'échanger avec vous lors d'un entretien.

\vspace{4pt}
Je vous prie d'agréer, Madame, Monsieur, l'expression de mes salutations distinguées.

\vspace{10pt}
\hfill Talvin Ackbaraly\hspace*{1cm}

\end{document}
